\section{Exercise 5}
\textit{Let $F$ be a vector field on $\mathbb{R}^3$ give by}
\[F(x,y,z) =\frac{1}{3}(x^2,y^2,z^3)\]

\textit{And let $C$ be the cube: $Cc= {(x,y,z) \in \mathbb{R}^3 | 0 \le x \le 1, 0 \le y \le 1, 0 \le z \le 1 }$}.

\begin{itemize}
    \item Calculate the flux out of the cube C directly
    \item Calculate the flux out of the cube C using Gauss Theorem (Divergence theorem)
\end{itemize}

Flux out of a cube, is the sum of the flux through all the flat surfaces of the cube.
And can be calculated as \[\Phi = \iint_S \mathbf{F} \cdot \mathbf{n} \ dS\]
The important detail here is that each surface of the cube, all six, must be calculated individually 
and then summed together.\\
Any of the six surfaces can be described by four points on any of the given two planes,
xy, xz, yz. For the ZY-plane, it will be bounded by four points: 
\[A = (0,0,0), \qquad B = (0,1,0), \qquad C = (0,1,1), \qquad D = (0,0,1)\]
This plane, which can also be called $x=0$ which is a side of the cube, is then the plane, which the flux will flow \textit{outwards} from,
meaning the normal vector to this plane must in the opposite direction of the cubes \textit{depth}. This leads to the normal vector:
\[N_{x0} = (-1,0,0)\] 
This rationality can then be applied to the rest of the cube as well.
\[
N_{x1} = (1,0,0), \qquad
N_{y0} = (0,-1,0), \qquad
N_{y1} = (0,1,0), \qquad
N_{z0} = (0,0,-1), \qquad
N_{z1} = (0,0,1), \qquad
\] 

The flux through a single face(or surface) of the cube, is the bounded between 0 and 1 in any 
direction, and because one point is fixed in place and we make minute changes in direction in the two others,
in the case of $x=0$, $dS = dy \ dz$ describes the small steps we take calculate the flux. 
So the flux becomes:

\[\Phi_{x=0} = \int_{0}^{1}\int_{0}^{1} (0,y,z) \cdot (-1,0,0)\ dy \ dz \Rightarrow\frac{1}{3}(0, y^2, z^2) \cdot (-1,0,0) \Rightarrow \frac{1}{3} \left(0 \cdot (-1) + y^2\cdot0 + z^2\cdot 0\right) = 0  \]
\[\Phi_{x=1} = \int_{0}^{1}\int_{0}^{1} (1,y,z) \cdot (1,0,0)\ dy \ dz \Rightarrow\frac{1}{3}(1, y^2, z^2) \cdot (1,0,0) \Rightarrow \frac{1}{3} \left(1 \cdot 1 + y^2\cdot0 + z^2\cdot 0\right) 
\]\[\Rightarrow \frac{1}{3} \int_{0}^{1}\int_{0}^{1}dy \ dz \Rightarrow \frac{1}{3} \int_{0}^{1} \left[y\right]_0^1  \Rightarrow \frac{1}{3} \int_{0}^{1} 1 dz \Rightarrow \frac{1}{3} \left[z\right]_0^1 = \frac{1}{3} \]

\[\Phi_{y=0} = \int_{0}^{1}\int_{0}^{1} (x,0,z) \cdot (0,-1,0)\ dx \ dz \Rightarrow\frac{1}{3}(x^2, 0, z^2) \cdot (0,-1,0) \Rightarrow \frac{1}{3} \left(x^2 \cdot 0 + 0 \cdot (-1) + z^2\cdot 0\right) = 0  \]
\[\Phi_{y=1} = \int_{0}^{1}\int_{0}^{1} (x,1,z) \cdot (0,1,0)\ dx \ dz \Rightarrow\frac{1}{3}(x^2, 1, z^2) \cdot (0,1,0) \Rightarrow \frac{1}{3} \left(x^2 \cdot 0 + 1\cdot 1 + z^2\cdot 0\right) = \frac{1}{3}  
\]\[\Rightarrow \frac{1}{3} \int_{0}^{1}\int_{0}^{1}dx \ dz \Rightarrow \frac{1}{3} \int_{0}^{1} \left[x\right]_0^1 \Rightarrow \frac{1}{3} \int_{0}^{1} 1 dz \Rightarrow \frac{1}{3} \left[z\right]_0^1 = \frac{1}{3} \]

\[\Phi_{z=0} = \int_{0}^{1}\int_{0}^{1} (x,y,0) \cdot (0,0,-1)\ dx \ dy \Rightarrow\frac{1}{3}(x^2, y^2, 0) \cdot (0,0,-1) \Rightarrow \frac{1}{3} \left(x^2 \cdot 0 + y^2 \cdot 0 + 0\cdot (-1)\right) = 0  \]
\[\Phi_{z=1} = \int_{0}^{1}\int_{0}^{1} (x,y,1) \cdot (0,0,1)\ dx \ dy \Rightarrow\frac{1}{3}(x^2, y^2, 1) \cdot (0,0,1) \Rightarrow \frac{1}{3} \left(x^2 \cdot 0 + y^2\cdot 0 + 1\cdot 1\right) = \frac{1}{3} 
\]\[\Rightarrow \frac{1}{3} \int_{0}^{1}\int_{0}^{1}dx \ dy \Rightarrow \frac{1}{3} \int_{0}^{1} \left[x\right]_0^1 \Rightarrow \frac{1}{3} \int_{0}^{1} 1 dy \Rightarrow \frac{1}{3} \left[y\right]_0^1 = \frac{1}{3} \]
This yields the final calculation for the flux out of the cube as:
\[\Phi_C =\Phi_{x=0} + \Phi_{x=1} + \Phi_{y=0} + \Phi_{y=1} + \Phi_{z=0} + \Phi_{z=1} = \frac{1}{3} + \frac{1}{3} + \frac{1}{3} = 1 \]